\subsection{Design Discussion and Chapter Summary}

This section summarizes the OTFS-IM design developed in this chapter and clarifies how it will be evaluated in Chapter 5. The purpose is not to introduce another processing block, but to connect the design choices into one coherent interpretation: OTFS-IM changes the information mapping in the delay-Doppler domain, while OHD and B-OHD change the geometry of the activation-pattern table used by that mapping.

\subsubsection{Summary of the OTFS-IM Design Chain}

The implemented OTFS-IM system is built as an extension of the baseline OTFS transceiver described in Chapter 3. The baseline OTFS modulation and demodulation blocks are reused. The main modification is introduced before OTFS modulation, where the full delay-Doppler QAM grid is replaced by a sparse index-modulated grid. In each IM block, \(b_1\) bits select an activation pattern and \(b_2\) bits select the QAM symbols transmitted on the active positions.

This structure gives OTFS-IM two information-bearing mechanisms. The first one is conventional symbol modulation, represented by the QAM symbols placed on active delay-Doppler positions. The second one is index modulation, represented by the locations of those active positions. Therefore, the receiver cannot simply demodulate QAM symbols independently. It must also infer which activation pattern was transmitted in each block. This is why the MP detector is used to produce soft probability information, and a block-wise MAP stage is then applied to jointly recover index bits and symbol bits.

Table~\ref{tab:chapter4_design_summary} summarizes the main design elements introduced in this chapter and their roles in the implemented simulation model.

\begin{table}[H]
    \centering
    \caption{Summary of the implemented OTFS-IM design elements}
    \label{tab:chapter4_design_summary}
    \begin{tabularx}{\linewidth}{>{\RaggedRight\arraybackslash}p{3.2cm}
                                  >{\RaggedRight\arraybackslash}X
                                  >{\RaggedRight\arraybackslash}p{3.2cm}}
        \toprule
        Design element & Role in the OTFS-IM system & Main evaluation aspect \\
        \midrule
        IM block \((n,k)\) &
        Defines how many positions are available and how many of them are active in each block. &
        Spectral efficiency and sparsity level \\
        \midrule
        Activation-pattern table \(\mathcal{P}\) &
        Maps index bits to valid activation patterns and constrains the receiver MAP search. &
        Pattern error rate and index BER \\
        \midrule
        Power normalization \(\alpha\) &
        Compensates for the fact that only \(k\) out of \(n\) positions carry active QAM symbols. &
        Fairness of BER comparison \\
        \midrule
        MP probability output &
        Provides soft probabilities for QAM symbols and the inactive zero state. &
        Reliability of block-wise decisions \\
        \midrule
        Block-wise MAP decision &
        Recovers the most likely activation pattern and active QAM symbols per IM block. &
        Total BER, index BER, symbol BER \\
        \midrule
        OHD and B-OHD selection &
        Select deterministic pattern tables based on Hamming-distance geometry and balance criteria. &
        Pattern reliability and random baseline validation \\
        \bottomrule
    \end{tabularx}
\end{table}

The most important point from Table~\ref{tab:chapter4_design_summary} is that OTFS-IM is not only a modulation-rate modification. It changes the error structure of the receiver. A wrong activation pattern produces index-bit errors and may also place detected QAM symbols in the wrong positions. For this reason, the simulation study must report more than total BER. The separate index BER, symbol BER, and pattern error rate are necessary to explain why one configuration performs better or worse than another.

\subsubsection{BER-SE Trade-Off and Pattern Reliability}

The use of index modulation creates a trade-off between spectral efficiency and reliability. For one IM block, the number of transmitted bits is \(b_1+b_2\), while the block occupies \(n\) delay-Doppler positions. Thus, the spectral efficiency of the OTFS-IM mapping is controlled by the block parameters \((n,k)\):
\begin{equation}
    \eta_{\mathrm{IM}}
    =
    \frac{b_1+b_2}{n}.
    \label{eq:chapter4_im_se_summary}
\end{equation}
Increasing the number of active positions generally increases the number of QAM symbol bits, while changing \(n\) and \(k\) also changes the number of index bits. However, a higher spectral efficiency does not automatically imply a better system. If the chosen patterns are difficult to distinguish, the additional index bits may increase the error rate.

This is the reason OHD and B-OHD are included in the design. Their purpose is to make the activation patterns more distinguishable. The OHD method focuses on the minimum Hamming distance between selected patterns, while B-OHD additionally considers average distance, the number of closest pairs, and carrier-usage balance. These criteria are not expected to remove all detection errors, especially at low SNR where the MP detector itself may produce uncertain probability estimates. Instead, they aim to reduce pattern confusion when the receiver has enough information to benefit from a better pattern geometry.

Therefore, the final evaluation should be read in two layers. The first layer compares baseline OFDM-LMMSE, baseline OTFS, and OTFS-IM to show the effect of moving from conventional multicarrier transmission to delay-Doppler modulation and then to index modulation. The second layer compares OHD, B-OHD, and the random-pattern diagnostic baseline to show whether deterministic pattern selection improves the reliability of OTFS-IM. If a random pattern table performs close to or better than OHD in some configuration, that result should not be hidden. It indicates that minimum-distance selection alone is not always sufficient and that the benefit of OHD depends on the selected \((n,k)\) setting.

\subsubsection{Quantities to Track in the OTFS-IM Evaluation}

Because OTFS-IM has both symbol-domain and index-domain information, the simulation must track several quantities at the same time. These quantities are not interchangeable. Spectral efficiency explains how many bits are transmitted per delay-Doppler resource element, total BER describes end-to-end bit reliability, and pattern error rate describes whether the receiver selected the correct activation pattern. If only one of these quantities is reported, the behavior of the system can be misinterpreted.

\begin{table}[H]
\centering
\caption{Main quantities used to interpret OTFS-IM results.}
\label{tab:otfs_im_evaluation_quantities}
\begin{tabularx}{0.95\linewidth}{|p{3.0cm}|X|}
\hline
\textbf{Quantity} & \textbf{Reason for reporting} \\
\hline
\(\eta_{\mathrm{IM}}\) & Shows whether an OTFS-IM configuration is rate-fair with baseline OTFS or belongs to a lower-rate reliability-oriented setting. \\
\hline
\(\mathrm{BER}_{\mathrm{IM}}\) & Measures the final end-to-end bit error probability after both index and symbol bits are reconstructed. \\
\hline
\(\mathrm{BER}_{\mathrm{idx}}\) & Isolates errors caused by incorrect recovery of the activation-pattern label. \\
\hline
\(\mathrm{BER}_{\mathrm{sym}}\) & Isolates errors in the QAM symbol bits on positions detected as active. \\
\hline
\(\mathrm{PER}_{\mathrm{pat}}\) & Measures block-level activation-pattern confusion, which is directly affected by pattern-table geometry. \\
\hline
\end{tabularx}
\end{table}

Table~\ref{tab:otfs_im_evaluation_quantities} is also useful for avoiding an overly simple conclusion. A configuration may have low total BER because it transmits fewer bits per block. Another configuration may have higher BER but preserve the same spectral efficiency as baseline OTFS. Similarly, a pattern table may reduce pattern error rate while the total BER remains limited by QAM symbol errors. The evaluation in Chapter 5 is therefore designed to connect these quantities rather than judge the system from a single curve.

\subsubsection{Connection to Chapter 5 Simulation Results}

The design in this chapter directly determines the structure of the simulation results in Chapter 5. The BER overview figure compares OFDM-LMMSE, baseline OTFS, OHD OTFS-IM, and B-OHD OTFS-IM under the same channel and SNR settings. The OTFS-IM error-decomposition figure then separates total BER into index BER, symbol BER, and pattern error rate so that the source of the errors can be identified. The pattern-geometry figure connects the selected activation table to Hamming-distance and carrier-usage properties. Finally, the random-pattern validation figure checks whether the proposed deterministic tables are actually better than arbitrary valid pattern selections.

This evaluation structure follows directly from the design logic of the chapter. If OTFS-IM improves BER at a comparable or higher spectral efficiency, the result supports the usefulness of index modulation in the delay-Doppler domain. If OHD or B-OHD reduces pattern error rate or index BER compared with random selection, the result supports the importance of pattern geometry. If the improvement is small or configuration-dependent, the result still has value because it identifies the limitation of the selected pattern design and motivates further work on detector-aware or channel-aware pattern selection.

In summary, this chapter has developed the OTFS-IM system used in the thesis simulations. It defined the IM block structure, spectral-efficiency expression, activation-pattern mapping, transmitter and receiver processing chains, block-wise MAP recovery, and the OHD/B-OHD pattern-selection methods. These elements provide the technical foundation for Chapter 5, where the proposed OTFS-IM design is evaluated through BER performance, error decomposition, pattern reliability, and BER-SE trade-off analysis.

\newpage
